\section{Curvature conditions for mass extrema}
\label{sec:mass-extremum-geometry}

At a smooth local extremum of the mass, the first two variations give a Gauduchon condition and a sign for one lower background power. On a threefold these necessary conditions are sufficient because the mass polynomial stops at its quadratic term. Fix a Hermitian form $\omega$ and use $\HH_\omega$ and $\Mass_\omega$ from \cref{chap:mass}.

For $n\geq2$, \cref{thm:mass-invariance-characterization,prop:gaud-volume-density-independence} identify the differential conditions for $\mathcal V_\mu(\omega)$ to be independent of the smooth positive probability measure $\mu$. By \cref{prop:entropy-mass-extrema}, these are also the conditions for $H_0-\operatorname{Ent}_{\mathrm{d}V}$ to be constant on finite-entropy probabilities, for a fixed smooth positive probability volume $\mathrm{d}V$; testing smooth positive probabilities suffices. In dimension one the mass is always $\int_X\omega$, so both conclusions hold automatically.

For the next argument, a smooth real $(n-1,n-1)$-form $T$ is nonnegative when $\mathrm i\alpha\wedge\bar\alpha\wedge T\geq0$ for every $(1,0)$-covector $\alpha$. This is the usual positivity in this bidegree. We need an elementary local test for its sign.

\begin{lemma}\label{lem:gradient-test-positive-form}
Suppose a smooth real $(n-1,n-1)$-form $T$ satisfies
\[
 \int_X \mathrm{d}f\wedge\dc f\wedge T\geq0
 \quad\hbox{for every smooth real }f.
\]
Then $T\geq0$.
\end{lemma}
\begin{proof}
If the conclusion fails at a point, choose a real linear function $\ell$ in local real coordinates, whose differential at that point has the offending complex direction. Shrinking the coordinate ball makes
$\mathrm{d}\ell\wedge\dc\ell\wedge T$ strictly negative there. Choose a nonzero smooth cutoff $\chi$ supported in this ball and, for positive integers $k$, set
\[
 f_k=k^{-1}\chi\sin(k\ell),
\]
extended by zero to $X$. Since $\mathrm{d}\ell$ is a nonzero constant covector in these coordinates,
\[
 \mathrm{d}f_k\wedge\dc f_k
 =\chi^2\cos^2(k\ell)\,\mathrm{d}\ell\wedge\dc\ell+O(k^{-1}),
\]
where the remainder has uniformly bounded coefficients after multiplication by $k$. The oscillatory part of $\cos^2(k\ell)=(1+\cos(2k\ell))/2$ has integral tending to zero, by integration by parts in a coordinate direction where $\ell$ has nonzero derivative. Hence
\[
 \lim_{k\to\infty}\int_X\mathrm{d}f_k\wedge\dc f_k\wedge T
 =\frac12\int_X\chi^2\mathrm{d}\ell\wedge\dc\ell\wedge T<0,
\]
a contradiction. Every complex covector is the $(1,0)$ part of a real covector, so this tests all required directions.
\end{proof}

\Needspace{12\baselineskip}
\begin{proposition}\label{prop:local-mass-extremum-necessary}
Let $n\geq3$, and let $v\in\HH_\omega$ be a local minimum of $\Mass_\omega$ in the $C^2$ topology. Then $\Omega=\omega+\ddc v$ satisfies
\[
 \ddc(\Omega^{n-1})=0,
 \qquad \ddc(\Omega^{n-2})\leq0.
\]
At a local maximum the first condition is unchanged and the second inequality is reversed.
\end{proposition}
\begin{proof}
Every smooth direction is admissible for a sufficiently short two-sided interval. The first variation gives the Gauduchon condition. Differentiating a second time gives
\[
 \begin{aligned}
 \left.\frac{\mathrm{d}^2}{\mathrm{d}t^2}\right|_{t=0}\Mass_\omega(v+tf)
 &=n(n-1)\int_X(\ddc f)^2\wedge\Omega^{n-2}\\
 &=-n(n-1)\int_X\mathrm{d}f\wedge\dc f\wedge\ddc(\Omega^{n-2}).
 \end{aligned}
\]
For the last equality, integrate once against $f\ddc f$, then use that $\ddc(\Omega^{n-2})$ is $\mathrm{d}$- and $\dc$-closed. At a minimum this second variation is nonnegative for every $f$. Apply \cref{lem:gradient-test-positive-form} to $-\ddc(\Omega^{n-2})$. The maximum case is identical with the opposite sign.
\end{proof}

In dimension three these necessary conditions are also sufficient, and a local extremum is automatically global.

\begin{theorem}\label{thm:threefold-mass-minimizer}
Assume $n=3$ and $v\in\HH_\omega$. The following are equivalent:
\begin{enumerate}
\item $v$ is a local minimum of $\Mass_\omega$;
\item $v$ is a global minimum of $\Mass_\omega$ on $\HH_\omega$;
\item $\Omega=\omega+\ddc v$ is Gauduchon and $\ddc\omega\leq0$.
\end{enumerate}
In this case $v_-(\omega)=\int_X\Omega^3$, and the Gauduchon infimum in \eqref{eq:gaud-lower-mass-formula} is attained. The corresponding statements for maxima hold with $\ddc\omega\geq0$, and then $v_+(\omega)=\int_X\Omega^3<+\infty$.
\end{theorem}
\begin{proof}
The preceding proposition gives the necessary conditions, since $\ddc\Omega=\ddc\omega$. Conversely, if $\Omega$ is Gauduchon, the dimension-three mass polynomial becomes, for every smooth real $f$,
\begin{equation}\label{eq:threefold-mass-dirichlet}
 \int_X(\Omega+\ddc f)^3-\int_X\Omega^3
       =-3\int_X\mathrm{d}f\wedge\dc f\wedge\ddc\omega.
\end{equation}
The form $\mathrm{d}f\wedge\dc f$ is nonnegative. The displayed difference is therefore nonnegative if $\ddc\omega\leq0$, proving a global minimum. Global minimality implies local minimality. Exact-background regularization identifies the infimum over smooth strictly positive potentials with $v_-(\omega)$, and \cref{prop:lower-mass-smooth-attainment} proves attainment in the Gauduchon formula. Reversing the sign in the same identity proves the maximum statements, including the equality with the supremum over bounded potentials.
\end{proof}

The local test supplies the necessity for smooth attainment in \eqref{eq:threefold-mass-dirichlet}. Neither that test nor the sufficient lower-bound argument produces such a representative in an arbitrary background.

\Needspace{8\baselineskip}
For the Hopf metric $\omega$ and the pulled-back Fubini--Study form $\alpha$ of \cref{chap:hopf}, $\ddc\omega=-\alpha^2\leq0$ and $\omega$ is Gauduchon. Thus the minimum is attained by the constant potential and equals $6$, consistently with its infinite upper mass. On a surface, by contrast, the mass is affine in the potential; a local extremum exists precisely when $\omega$ is Gauduchon, in which case all masses are equal. In higher dimensions the necessary sign condition of \cref{prop:local-mass-extremum-necessary} does not control all higher terms of the mass polynomial, and sufficiency is not claimed.
